% !TEX root = ../../../main.tex
% Block card: preprocessing.highpass -- "Highpass filter" (Part I, the blocks chapter).
\cardfield{What it does}
Removes the slow drift from a signal and keeps the faster movements riding on it.

\cardfield{When to use it}
Electrodes in a fungal culture drift: the level wanders over hours while the activity you care
about comes and goes over seconds to minutes. Picture a boat's deck rising and falling on a
long swell while someone jumps on it: a high-pass filter takes away the swell and leaves the
jumps. Use it before a detector whose threshold assumes the signal sits near zero. If big drops
or spikes dominate the recording, \blockref{preprocessing.beads} removes drift with less damage
to them.

\cardfield{What goes in, what comes out}
In goes a span of one channel. Out comes a signal of the same length and sampling rate, centred
on zero, with everything slower than the cut-off taken away.

\cardfield{Settings that matter}
\emph{cutoff\_hz} is the dividing speed, in hertz. The default, 0.01\,Hz, removes changes
that take longer than about 100 seconds; lower it to keep slower events. \emph{order} sets
how sharply the cut is made.

\cardfield{Watch out for}
An event slower than the cut-off is removed along with the drift: a drop that recovers over
ten minutes is mostly flattened by the default. A sharp drop gains a small rise before it and
an overshoot after it. One missing sample turns the whole output into missing values, and the
cut-off must be below half the sampling rate.
